\documentclass[../Main.tex]{subfiles}
\begin{document}
\chapter{2014-2015学年微积分（一）（上）期中考试}

\section{基本计算题(每小题 6 分，共 60 分)}
\begin{enumerate}
    \item 计算数列极限$l=\lim_{n\to\infty}\left(1+n^{2}+2^{n}\right)^{\frac{1}{n}}$.
          \vspace{11em}

    \item 计算极限$l=\lim_{x\to 0}\frac{\arctan x-\sin x}{\ln(1+x^{3})}$.
          \vspace{11em}

    \item 计算极限$l=\lim_{x\to 0}\left(\frac{\tan x}{x}\right)^{\frac{1}{1-\cos
                      x}}$.
          \vspace{11em}

    \item 已知$y=f\left(\frac{x-1}{x+1}\right)$，其中$f'(x)=\arcsin x^{2}$，求$\left
              .\frac{\mathrm{d}y}{\mathrm{d}x}\right|_{x=0}$.
          \vspace{9em}

    \item 设参数方程$\begin{cases}
                  x=t-\ln(1+t), \\
                  y=t^{3}+t^{2}
              \end{cases}$确定了函数$y=y(x)$，求$\frac{\mathrm{d}y}{\mathrm{d}x}$，$\frac{\mathrm{d}^{2}
                  y}{\mathrm{d}x^{2}}$.
          \vspace{9em}

    \item 计算曲线$x^{2}-xy+2y^{2}=2$在点$(1,1)$处的切线方程.
          \vspace{9em}

    \item 设$x=g(y)$是函数$y=\ln x+\arctan x$的反函数，求$y=\frac{\pi}{4}$处的导数$\frac{\mathrm{d}x}{\mathrm{d}y}$和$\frac{\mathrm{d}^{2}
                  x}{\mathrm{d}y^{2}}$.
          \vspace{9em}

    \item 设函数$f(x)$二阶可导，计算下列函数的导数$y'$以及$y''$：

          （1）$y=f(x^{2})$；

          （2）$y=(f(x))^{2}$.
          \vspace{9em}

    \item 设函数$y=\frac{(1+x)^{2}\sqrt{x}}{x^{5}\mathrm{e}^{x}}$，使用对数求导法求$\left
              .y'\right|_{x=1}$.
          \vspace{11em}

    \item 求无穷小量$u(x)=\cos 2x-\frac{1}{\mathrm{e}^{2x^2}}(x\to0)$的主部.
          \vspace{11em}
\end{enumerate}

\section{综合题(每小题 7 分，共 28 分)}
\begin{enumerate}
    \setcounter{enumi}{10}
    \item 设函数$f(x)=
              \begin{cases}
                  \frac{\cos x}{x+2},            & x\geq0, \\
                  \frac{\sqrt{a}-\sqrt{a-x}}{x}, & x<0
              \end{cases}$，问$a(a\geq0)$为何值时，$x=0$是$f(x)$的间断点，并指出该间断点的类型.
          \vspace{12em}

    \item 设函数$f(x)$在$x=0$处二阶可导.$f'(0)=0,f''(0)=2$.求$l=\lim_{x\to0}
              \frac{f(\tan x)-f(x)}{x^{4}}$.
          \vspace{12em}

    \item 设$f'(x)$处处连续，$g(x)=f(x)\sin^{2} x$，求$g''(0)$.
          \vspace{12em}

    \item 一个13英尺长的梯子斜靠在墙边上，当梯子的顶端沿着墙面向墙底滑落时，梯子底端沿地面移动的速度是5英尺/秒，问当梯子底端的地面长度为12英尺时，直角三角形的面积的变化率是多少？
          \vspace{11em}
\end{enumerate}

\section{证明题(每小题 6 分，共 12 分)}
\begin{enumerate}
    \setcounter{enumi}{14}
    \item 设函数$f(x),g(x)$在$[a,b]$上连续，在$(a,b)$内可导，证明存在$\xi\in
              (a,b)$，使得$\frac{f'(\xi)}{g'(\xi)}=\frac{f(\xi)-f(a)}{g(b)-g(\xi)}$.
          \vspace{12em}

    \item 设函数$f(x)=\alpha_{1}\varphi(x)+\alpha_{2}\varphi(2x)+\cdots+\alpha
              _{n}\varphi(nx)$，其中$\alpha_{1},\alpha_{2},\cdots,\alpha_{n}$是常数，$\varphi
              (0)=0$，$\varphi'(0)=1$，已知对一切实数$x$，有$|f(x)|\leq |x|$，试证：$|\alpha
              _{1}+2\alpha_{2}+\cdots+n\alpha_{n}|\leq 1$.
          \vspace{12em}
\end{enumerate}

\chapter{2014-2015学年微积分（一）（上）期中考试参考答案}
\section{基本计算题(每小题 6 分，共 60 分)}
\begin{enumerate}
    \item \textbf{Solution}. 当$n$充分大时，$2^{n}>1+n^{2}$，所以$2^{n}<1+n^{2}+2
                  ^{n}<2\cdot 2^{n}$，因此
          \[
              2<\left(1+n^{2}+2^{n}\right)^{\frac{1}{n}}<2^{1+\frac{1}{n}}.
          \]

          又$\lim_{n\to\infty}2^{1+\frac{1}{n}}=2$，

          由夹逼定理得$\lim_{n\to\infty}\left(1+n^{2}+2^{n}\right)^{\frac{1}{n}}=2$.

    \item \textbf{Solution}.
          \[
              \begin{aligned}
                  l = & \lim_{x\to 0}\frac{\arctan x-\sin x}{x^3}                                                          \\
                  =   & \lim_{x\to0}\frac{\frac{1}{1+x^2}-\cos x}{3x^2}=\lim_{x\to0}\frac{1-\cos x-x^2\cos x}{3x^2(1+x^2)} \\
                  =   & \lim_{x\to0}\frac{1-\cos x}{3x^2}- \lim_{x\to0}\frac{x^2\cos x}{3x^2}                              \\
                  =   & \frac{1}{6}-\frac{1}{3}=-\frac{1}{6}.
              \end{aligned}
          \]

    \item \textbf{Solution}.
          \[
              \begin{aligned}
                  l= & \exp\left(\lim_{x\to0}\frac{\ln\tan x-\ln x}{1-\cos x}\right)                                                                                 \\
                  =  & \exp\left(\lim_{x\to0}\frac{\frac{\sec^2 x}{\tan x}-\frac{1}{x}}{\sin x}\right)=\exp\left(\lim_{x\to0}\frac{x-\sin x\cos x}{x\sin^2 x}\right) \\
                  =  & \exp\left(\lim_{x\to0}\frac{1-\cos 2x}{3x^2}\right)                                                                                           \\
                  =  & \mathrm{e}^{\frac{2}{3}}.                                                                                                                     \\
              \end{aligned}
          \]

    \item \textbf{Solution}. $y'=f'\left(\frac{x-1}{x+1}\right)\cdot\frac{2}{(x+1)^{2}}$，所以$\left
              .\frac{\mathrm{d}y}{\mathrm{d}x}\right|_{x=0}=f'(-1)\cdot 2 =2\arcsin 1=\pi$.

    \item \textbf{Solution}. $\frac{\mathrm{d}y}{\mathrm{d}x}=\frac{\frac{\mathrm{d}y}{\mathrm{d}t}}{\frac{\mathrm{d}x}{\mathrm{d}t}}
              =\frac{3t^{2}+2t}{1-\frac{1}{1+t}}=(1+t)(3t+2)$，

          $\frac{\mathrm{d}^{2}y}{\mathrm{d}x^{2}}=\frac{\mathrm{d}}{\mathrm{d}t}\left
              (\frac{\mathrm{d}y}{\mathrm{d}x}\right)\cdot\frac{\mathrm{d}t}{\mathrm{d}x}
              =\left((1+t)(3t+2)\right)'\cdot\frac{1}{1-\frac{1}{1+t}}=\frac{(6t+5)(1+t)}{t}$.

    \item \textbf{Solution}. 方程$x^{2}-xy+2y^{2}=2$两边求微分得
          \[
              2x\,\mathrm{d}x-(x\,\mathrm{d}y+y\,\mathrm{d}x)+4y\,\mathrm{d}y=0,
          \]

          将$x=1,y=1$代入上式得$\mathrm{d}y=-\frac{1}{3}\,\mathrm{d}x$，所以切线方程为$y
              -1=-\frac{1}{3}(x-1)$即$x+3y-4=0$.

    \item \textbf{Solution}. 函数$y=\ln x+\arctan x$严格单调增加，当$y=\frac{\pi}{4}$时，$x
              =1$.

          $\frac{\mathrm{d}x}{\mathrm{d}y}=\frac{1}{\frac{\mathrm{d}y}{\mathrm{d}x}}=
              \frac{1}{\frac{1}{x}+\frac{1}{1+x^{2}}}=\frac{x(1+x^{2})}{x^{2}+x+1}$，所以$\left
              .\frac{\mathrm{d}x}{\mathrm{d}y}\right|_{y=\frac{\pi}{4}}=\frac{1\cdot 2}{3}
              =\frac{2}{3}$.

          \[
              \begin{aligned}
                  \frac{\mathrm{d}^{2} x}{\mathrm{d}y^{2}}= & \frac{\mathrm{d}}{\mathrm{d}x}\left(\frac{\mathrm{d}x}{\mathrm{d}y}\right)\cdot\frac{\mathrm{d}x}{\mathrm{d}y} \\
                  =                                         & \frac{\mathrm{d}}{\mathrm{d}x}\left(\frac{x(1+x^{2})}{x^{2}+x+1}\right)\cdot\frac{x(1+x^{2})}{x^{2}+x+1}       \\
                  =                                         & \frac{(3x^{2}+1)(x^{2}+x+1)-x(1+x^{2})(2x+1)}{(x^{2}+x+1)^{2}}\cdot\frac{x(1+x^{2})}{x^{2}+x+1}.
              \end{aligned}
          \]

          所以$\left.\frac{\mathrm{d}^{2} x}{\mathrm{d}y^{2}}\right|_{y=\frac{\pi}{4}}
              =\frac{2}{3}\cdot\frac{4\cdot 3-2\cdot 3}{3^{2}}=\frac{4}{9}$.

    \item \textbf{Solution}. （1） $y'=f'(x^{2})\cdot 2x$，$y''=f''(x^{2})\cdot (
              2x)^{2}+f'(x^{2})\cdot 2 = 4x^{2} f''(x^{2})+2f'(x^{2})$.

          （2）$y'=2f(x)f'(x)$，$y''=2(f'(x))^{2}+2f(x)f''(x)$.

    \item \textbf{Solution}. $\ln y = 2\ln(1+x)+\frac{1}{2}\ln x-5\ln x-x$，所以
          \[
              y' = \frac{(1+x)^{2}\sqrt{x}}{x^{5}\mathrm{e}^{x}}\left(\frac{2}{x+1}-\frac{9}{2x}
              -1\right).
          \]

          $\left.y'\right|_{x=1}=\frac{4}{\mathrm{e}}\left(1-\frac{9}{2}-1\right)=-\frac{18}{\mathrm{e}}$.

    \item \textbf{Solution}.
          \[
              \begin{aligned}
                  u(x) = & \cos x-\mathrm{e}^{-2x^2}                                                                                          \\
                  =      & \left(1-\frac{1}{2}(2x)^{2}+\frac{1}{24}(2x)^{4}+o(x^{4})\right)-\left(1-2x^{2}+\frac{(2x^2)^2}{2}+o(x^{4})\right) \\
                  =      & -\frac{4}{3}x^{4}+o(x^{4}).
              \end{aligned}
          \]

          所以无穷小量$u(x)$的主部为$-\frac{4}{3}x^{4}$.
\end{enumerate}

\section{综合题(每小题 7 分，共 28 分)}
\begin{enumerate}
    \setcounter{enumi}{10}
    \item \textbf{Solution}. $\lim_{x\to 0^+}f(x)=\lim_{x\to 0^+}\frac{\cos
                  x}{x+2}=\frac{1}{2}$.

          若$a=0$，则$\lim_{x\to 0^-}f(x)=\lim_{x\to 0^-}\frac{-\sqrt{-x}}{x}=+\infty$，此时$x
              =0$为无穷间断点.

          若$a\neq0$，
          \[
              \lim_{x\to 0^-}f(x)=\lim_{x\to 0^-}\frac{\sqrt{a}-\sqrt{a-x}}{x}=\lim_{x\to
                  0^-}\frac{(\sqrt{a}-\sqrt{a-x})(\sqrt{a}+\sqrt{a-x})}{x(\sqrt{a}+\sqrt{a-x})}
              =\lim_{x\to 0^-}\frac{x}{x(\sqrt{a}+\sqrt{a-x})}=\frac{1}{2\sqrt{a}}.
          \]

          若$a\neq 1$，$\lim_{x\to 0^+}f(x)=\frac{1}{2\sqrt{a}}\neq \frac{1}{2}$，此时$x
              =0$为跳跃间断点.

    \item \textbf{Solution}. 由于$f''(0)$存在，则$f(x)$在$x=0$的邻域内可导.

          由Lagrange中值定理，存在$\xi$介于$\tan x$与$x$之间，使得
          \[
              \begin{aligned}
                  \lim_{x\to0}\frac{f(\tan x)-f(x)}{x^4} & =\lim_{x\to0}\frac{f'(\xi)(\tan x-x)}{x^4}                                           \\
                                                         & =\lim_{x\to0}\frac{f'(\xi)-f'(0)}{\xi-0}\cdot\frac{\xi}{x}\cdot\frac{\tan x-x}{x^3}.
              \end{aligned}
          \]

          当$x\to0$时，$\tan x\to0$，所以由夹逼准则可知$\xi\to0$.

          同时，$\frac{\xi}{x}$介于$1$与$\frac{\tan x}{x}$之间，显然$\lim_{x\to0}\frac{\xi}{x}
              =1$. 故
          \[
              \begin{aligned}
                  \lim_{x\to0}\frac{f'(\xi)-f'(0)}{\xi-0}\cdot\frac{\xi}{x}\cdot\frac{\tan x-x}{x^3} & =f''(0)\cdot\frac{\sec^2 x-1}{3x^2} \\
                                                                                                     & =\frac{2}{3}.
              \end{aligned}
          \]

    \item \textbf{Solution}. $g'(x)=f'(x)\sin^{2} x+f(x)\cdot 2\sin x\cos x$，所以$g
              '(0)=0$.
          \[
              \begin{aligned}
                  g''(0) = & \lim_{x\to0}\frac{g'(x)-g'(0)}{x}                                                 \\
                  =        & \lim_{x\to0}\frac{f'(x)\sin^2 x+f(x)\cdot 2\sin x\cos x}{x}                       \\
                  =        & f'(0)\lim_{x\to0}\frac{\sin^2 x}{x}+\lim_{x\to 0}2f(x)\cos x\cdot\frac{\sin x}{x} \\
                  =        & 2f(0).
              \end{aligned}
          \]

    \item \textbf{Solution}. 设$t$时刻梯子底端的地面长度为$x(t)$英尺，梯子顶端距地面的距离为$y
              (t)$英尺，

          直角三角形的面积为$S(t)$平方英尺. 则$x^{2}+y^{2}=13^{2}$，$S=\frac{1}{2}xy$.

          方程$x^{2}+y^{2}=13^{2}$两边对$t$求导得$x\frac{\mathrm{d}x}{\mathrm{d}t}+y\frac{\mathrm{d}y}{\mathrm{d}t}
              =0$，

          此时刻$x=12$，所以$y=5$，又$\frac{\mathrm{d}x}{\mathrm{d}t}=5$，代入得$\frac{\mathrm{d}y}{\mathrm{d}t}
              =-\frac{12}{5}\cdot 5=-12$.

          故$\frac{\mathrm{d}S}{\mathrm{d}t}=\frac{1}{2}\left(x\frac{\mathrm{d}y}{\mathrm{d}t}
              +y\frac{\mathrm{d}x}{\mathrm{d}t}\right)=\frac{1}{2}(12\cdot(-12)+5\cdot 5)
              =-\frac{119}{2}$.

          即直角三角形的面积的变化率为$-\frac{119}{2}$平方英尺/秒.
\end{enumerate}

\section{证明题(每小题 6 分，共 12 分)}
\begin{enumerate}
    \setcounter{enumi}{14}
    \item \textbf{Proof}. 令$F(x)=f(x)g(x)-f(x)g(b)-f(a)g(x)$，则函数$F(x)$在$[
                  a,b]$上连续，在$(a,b)$内可导，

          且$F(a)=F(b)=-f(a)g(b)$.由Rolle定理，存在$\xi\in(a,b)$使得$F'(\xi)=0$，

          即$f'(\xi)g(\xi)+f(\xi)g'(\xi)-f'(\xi)g(b)-f(a)g'(\xi)=0$，化简得$\frac{f'(\xi)}{g'(\xi)}
              =\frac{f(\xi)-f(a)}{g(b)-g(\xi)}$.

    \item \textbf{Proof}. $f(0)=\alpha_{1}\varphi(0)+\alpha_{2}\varphi(0)+\cdots
              +\alpha_{n}\varphi(0)=0$.
          \[
              \begin{aligned}
                  f'(0)= & \lim_{x\to0}\frac{f(x)-f(0)}{x-0}=\lim_{x\to0}\frac{\alpha_1\varphi(x)+\alpha_2\varphi(2x)+\cdots+\alpha_n\varphi(nx)}{x}                                  \\
                  =      & \lim_{x\to 0}\alpha_{1}\frac{\varphi(x)}{x}+\lim_{x\to0}\alpha_{2}\frac{\varphi(2x)}{2x}\cdot 2+\cdots+\lim_{x\to0}\alpha_{n}\frac{\varphi(nx)}{nx}\cdot n \\
                  =      & \varphi'(0)\left(\alpha_{1}+2\alpha_{2}+\cdots+n\alpha_{n}\right)=\alpha_{1}+2\alpha_{2}+\cdots+n\alpha_{n}.
              \end{aligned}
          \]

          因为$\forall\,x$，$|f(x)|\leq |x|$，所以$\left|\frac{f(x)-f(0)}{x-0}\right|
              \leq 1$，令$x\to0$，得$|f'(0)|=|\alpha_{1}+2\alpha_{2}+\cdots+n\alpha_{n}|\leq
              1$.
\end{enumerate}
\end{document}
